% TOP_PROBLEM: {"id": 1282, "title": "Irrationality of Euler's Constant", "category_id": 1, "category": {"id": 1, "name": "number_theory", "display_name": "Number Theory", "description": "Properties of integers, prime numbers, Diophantine equations.", "slug": "number-theory", "order_index": 1, "created_at": "2026-07-31T15:26:25.670Z"}, "status": "open"}
% FIRST_POSTED: 2026-09-25
% ENTRY_KIND: statement_only
% !TeX program = lualatex
% Definition and sources only; this file is not an LLM solution attempt.
\documentclass[11pt]{article}
\usepackage[margin=25mm]{geometry}
\usepackage{fontspec}
\setmainfont{Latin Modern Roman}
\usepackage{amsmath,amssymb,mathtools}
\usepackage{unicode-math}
\setmathfont{Latin Modern Math}
\usepackage{xurl,hyperref}
\hypersetup{colorlinks=true,urlcolor=blue}
\setcounter{secnumdepth}{0}
\setlength{\parindent}{0pt}
\setlength{\parskip}{5pt}
\setlength{\emergencystretch}{3em}
\begin{document}
\section*{\texorpdfstring{Irrationality of Euler's Constant}{Irrationality of Euler's Constant}}\label{1282}
\subsection{Definitions and mathematical statement}
Euler's constant is defined using natural logarithms by
\[
 \gamma=\lim_{n\to\infty}\left(H_n-\log n\right),
 \qquad H_n=\sum_{j=1}^n\frac1j.
\]
The selected question is whether $\gamma\notin{\mathbb{Q}}$, equivalently whether
$q\gamma-p\ne0$ for every $p\in{\mathbb{Z}}$ and integer $q\ge1$. The sequence defining the limit, rather than a finite decimal expansion, specifies the constant exactly. This is not Euler's exponential constant $e$, and transcendence of $\gamma$ would be a stronger conclusion than the requested irrationality.
\par\smallskip\textit{Formulation and concept references:} \href{https://www.ams.org/journals/bull/2013-50-04/S0273-0979-2013-01423-X/S0273-0979-2013-01423-X.pdf}{[S1]}.

\subsection{Short English statement}
Is the limiting difference between the harmonic sum and the natural logarithm an irrational number?
\subsection{Sources}
\begin{itemize}
\item[S1]Jeffrey C. Lagarias, "Euler's constant: Euler's work and modern developments," Bull. Amer. Math. Soc. 50 (2013), no. 4, 527-628.
\url{https://www.ams.org/journals/bull/2013-50-04/S0273-0979-2013-01423-X/S0273-0979-2013-01423-X.pdf}.
\textit{Formulation source.}
\end{itemize}
\end{document}
